\scanpage{15}
\mspage{12}

\margin{$\times$}

Intuitivement, on regarde un dérivateur $\mathbb D$ comme ``étant'' la catégorie fondamentale $\unsure{\mathcal A|\mathscr A}$, avec une ``structure'' supplémentaire, savoir la donnée des ``extensions'' de $\mathcal A$ \margin{$\longrightarrow$}\inserted{(comme cat. de coeff. \uline{fond.})} aux diagrammes $I\in\mathrm{Diag}$, de façon contravariante en $I$. Un \deleted{objet}\inserted{coeff $F$} de $\deleted{\unsure{\mathcal A}}\mathbb D(I)\overset{\mathrm{d\acute ef}}{=}\mathcal A_{I}$ est regardé comme ``étant'' le diagramme $(F_{i})_{i\in I}$ correspondant d'objets fondamentaux, avec une donnée supplémentaire de ``globalisation''. Les $F_{i}$, \deleted{\unsure{sections}} \uline{ou} ``fibres'' de $F$ en les $i\in I$, sont les caractéristiques \unsure{locaux|locales}, du coéfficient ``global sur $I$'' $F$.

\figure{\begin{tikzpicture}\draw (0,.43) -- (.28,0);\end{tikzpicture}}

Tous les dérivateurs connus sont obtenus de la façon \unsure{standart|standard} suivante. On part d'une ``catégorie de modèles''

\[
\unsure{\mathcal M|\mathscr M},\ \Sigma
\]

munie d'un ens. de flèches, appelées ``équivalences'', et satisfaisant aux axiomes pertinents pour le passage à une catégorie de fractions,

\[
\mathcal A=\mathcal M\Sigma^{-1}\ \unsure{,}
\]

qui sera la ``catégorie fondamentale''. On définit pour \unsure{tt|tte} $I$ \unsure{dans|de} $\mathrm{Diag}$

\[
\begin{aligned}
\mathcal M_{I}&=\underline{\mathrm{Hom}}\,(I^{\unsure{\circ|o}},\mathcal M)\\
\Sigma_{I}&=\{u\in\mathrm{Fl}(\mathcal M_{I})\mid\forall i\in I,\ u_{i}\in\Sigma\}\\
\mathcal A_{I}&=\mathcal M_{I}\Sigma_{I}^{-1}
\end{aligned}
\]

% archive sheet number: 14
